\section{Introduction}
\label{sec:introduction}

Unmanned aerial vehicles (UAVs) have become an important component of flexible wireless communication systems. Compared with fixed terrestrial infrastructure, UAV communication platforms can be rapidly deployed, repositioned according to demand, and used for emergency coverage, temporary networking, surveillance support, aerial relaying, and mission-oriented connectivity \cite{zeng2016uav,mozaffari2019tutorial}. The same mobility and line-of-sight channel characteristics that make UAVs attractive, however, also create security risks. Wireless signals are broadcast by nature, UAV-ground links are often visible over large areas, and adversarial receivers can move to favorable interception locations. Physical-layer security is therefore a natural tool for UAV systems because it directly characterizes secrecy performance from the difference between legitimate and eavesdropping channels \cite{wyner1975wiretap,csiszar1978broadcast}.

A substantial body of work has studied secure UAV communication through trajectory optimization, power allocation, and cooperative jamming. These studies have shown that UAV mobility can be actively exploited to improve secrecy rate by moving the legitimate transmitter closer to intended users and by positioning a friendly jammer to degrade the eavesdropper channel \cite{zhang2019securing,zhong2019cooperative,zhou2018unknown,sun2019aerial}. Most such formulations, however, assume that the controller either knows the relevant channel state or has a fixed uncertainty model. In practical secure UAV operation, the controller may only maintain an estimated state of the eavesdropper through sensing, tracking, network telemetry, or edge-side prediction. This state estimate can be viewed as a digital twin of the adversarial environment.

Digital twins are increasingly discussed as enablers for autonomous wireless networks and UAV networks because they can support simulation, prediction, policy adaptation, and network management before actions are applied to the physical system \cite{kuruvatti2022digital,mcmanus2023digital}. A twin is useful only when it remains sufficiently synchronized with the physical process. Updating the twin consumes sensing, control, or communication resources; delayed or failed updates leave the controller with stale information. This creates a control problem that is different from conventional trajectory-only optimization: the system must decide not only where the UAVs should move and how much power they should use, but also when and how strongly the twin should be synchronized.

This paper studies this coupled problem for secure UAV communication. We consider two cooperative UAVs, multiple ground users, a mobile eavesdropper, and a limited synchronization budget. One UAV serves as the confidential transmitter, while the other provides cooperative jamming. The controller observes a digital-twin state of the eavesdropper and must choose synchronization bandwidth, UAV motion, source power, and jamming power at every time slot. The objective is to improve average secrecy rate, reduce secrecy outage, and avoid unsafe decisions caused by stale or uncertain twin states. Figure~\ref{fig:framework} summarizes the resulting closed-loop control structure.

\begin{figure*}[t]
\centering
\includegraphics[width=0.92\textwidth]{dt_uav_joint_sync_secure_control_loop.pdf}
\caption{Digital-twin-aware joint synchronization and secure control loop. The physical UAV system updates the digital-twin state through budget-limited sensing and synchronization. The risk-aware rollout controller uses twin features and secrecy-risk slack to jointly select synchronization, UAV motion, source-power, and jamming-power actions.}
\label{fig:framework}
\end{figure*}

The main challenge is that synchronization and communication control are strongly coupled over time. A fixed-period synchronization rule is simple and inexpensive, but it does not react to secrecy risk, remaining budget, or the future value of information. A security-risk trigger can synchronize under dangerous states, but it still separates synchronization from motion and power decisions. A local optimization method can improve per-slot actions, but it may be shortsighted when future twin quality and outage penalties matter. A learning-based controller can be fast at inference time, but the coupled action space makes training difficult under limited data and computation.

To address these issues, we propose an empirical-risk-aware joint rollout controller, denoted as \texttt{rollout\_joint}. The controller enumerates feasible synchronization bandwidths, UAV movements, source power levels, and jamming power levels, and evaluates them using a finite-horizon score. The score combines predicted secrecy reward, outage penalty, synchronization cost, twin-quality terms, and an empirical secrecy-loss penalty. The fitted risk estimator is learned from simulated trajectories using a holdout procedure and is used as an empirical upper bound on secrecy loss induced by twin uncertainty. We emphasize that this estimator is an in-policy, simulation-validated diagnostic; it is not claimed as a finite-sample safety guarantee or a universal proof under arbitrary distribution shift.

The contributions of this paper are as follows:
\begin{enumerate}
    \item We formulate a digital-twin-aware secure UAV communication problem in which twin synchronization, secrecy rate, secrecy outage, UAV mobility, transmit power, and cooperative jamming are optimized in a unified control loop.
    \item We design an empirical-risk-aware rollout controller that jointly selects synchronization and communication-control actions, rather than treating twin updates as a separate rule-based module.
    \item We construct a holdout-fitted secrecy-loss risk estimator and report empirical upper-bound coverage on held-out trajectories from base, hard, and calibrated stress scenarios.
    \item We provide a multi-seed simulation study with rule-based baselines, SCA baselines, a lightweight PPO diagnostic run, oracle references, ablations, runtime analysis, and a small exact-DP sanity check. In the calibrated stress scenario, the proposed controller improves average secrecy rate over periodic synchronization by 0.2683 and reduces secrecy outage by 0.2141.
\end{enumerate}

The rest of this paper is organized as follows. Section~\ref{sec:related} reviews related work. Section~\ref{sec:model} presents the system model and metrics. Section~\ref{sec:method} describes the risk estimator and joint rollout method. Section~\ref{sec:experiments} reports the experimental results. Section~\ref{sec:discussion} discusses interpretation and limitations. Section~\ref{sec:conclusion} concludes the paper.
